\begingroup\small\renewcommand{\arraystretch}{1.12}\setlength{\LTcapwidth}{\textwidth}
\setlength{\tabcolsep}{4pt}
\begin{longtable}{@{}lp{8.0cm}cl@{}}
\caption{All positively identified rings in the census. U means a unitary realization; S means a spherical realization and no pseudo-unitary categorification (hence no unitary one). The numbers of types U and S are given in Table~\ref{tab:classification}. Here $R_{m,j}=R^{(6)}_{m,j}$. Every constructor and its based isomorphism are recorded electronically.}\label{tab:models}\\
\toprule Ring & Model or construction & Type & Source\\\midrule\endfirsthead
\multicolumn{4}{l}{\tablename\ \thetable\ continued}\\
\toprule Ring & Model or construction & Type & Source\\\midrule\endhead
\midrule\multicolumn{4}{r}{Continued on the next page}\\\endfoot
\bottomrule\endlastfoot
$R_{1,1}$ & $\mathcal I(C_4,1)$: cyclic, non-self-dual Ising zest & U & \cite{LPR}\\
$R_{1,3}$ & $\mathcal I(C_4,0)$: cyclic, self-dual Ising zest & U & \cite{LPR}\\
$R_{1,4}$ & $\Rep(C_3\rtimes C_4)$ & U & \cite{LPR}\\
$R_{1,6}$ & $\Vec_{C_2}\boxtimes\Vec_{C_3}\simeq\Vec_{C_6}$ & U & \cite{LPR}\\
$R_{1,7}$ & Tambara--Yamagami category for $C_5$ & U & \cite{LPR}\\
$R_{1,9}$ & $\mathcal I(C_2^2,1)$: non-self-dual Ising zest & U & \cite{LPR}\\
$R_{1,10}$ & Zested $SO(5)_2$ category & U & \cite{LPR}\\
$R_{1,11}$ & $\mathrm{Fib}\boxtimes\Vec_{C_3}$ & U & \cite{LPR}\\
$R_{1,14}$ & $\Vec_{S_3}$ & U & \cite{LPR}\\
$R_{1,15}$ & Haagerup category $\mathcal H_3$ & U & \cite{GS}\\
$R_{1,20}$ & $\Vec_{C_2}\boxtimes\mathrm{Ising}$; also $R(0,0)$ & U & \cite{LPR,MY}\\
$R_{1,21}$ & $\Vec_{C_2}\boxtimes\Rep(S_3)$; also $R(0,1)$ & U & \cite{LPR,MY}\\
$R_{1,23}$ & $\Vec_{C_2}\boxtimes PSU(2)_5$ & U & \cite{LPR}\\
$R_{1,24}$ & $PSU(2)_{10}$ (the other $3311$ even-part ring) & U & \cite{LPR,IJMS}\\
$R_{1,26}$ & $SO(5)_2$ & U & \cite{LPR}\\
$R_{1,30}$ & $\Rep(D_9)$, with $|D_9|=18$ & U & \cite{LPR}\\
$R_{1,32}$ & $\Rep((C_3\times C_3)\rtimes C_2)$, inversion action & U & \cite{LPR}\\
$R_{1,34}$ & $PSU(2)_{11}$ & U & \cite{LPR}\\
$R_{1,36}$ & $\mathrm{Fib}\boxtimes\mathrm{Ising}$ & U & \cite{LPR}\\
$R_{1,37}$ & $\mathrm{Fib}\boxtimes\Rep(S_3)$ & U & \cite{LPR}\\
$R_{1,38}$ & $\mathrm{Fib}\boxtimes PSU(2)_5$ & U & \cite{LPR}\\
\addlinespace
$R_{2,2}$ & Fermionic $C_2$-zest of $\operatorname{Ad}(E_6)\boxtimes\Vec_{C_2}$ & U & \cite{DGP}\\
$R_{2,4}$ & $\mathcal E_4(G_2)$: modules for $G_{2,4}\subset\operatorname{Spin}(14)_1$ & U & \cite{CRT}\\
$R_{2,11}$ & $\Rep((C_3\times C_3)\rtimes C_4)$ & U & \cite{EGNO}\\
$R_{2,14}$ & Fermionic grading zest of the $E_6$ category & U & \cite{DGP}\\
$R_{2,34}$ & $TT_3$ & U & \cite{EM}\\
$R_{2,36}$ & $\mathcal C(A_4,1,C_2,1)$ & U & \cite{Integral}\\
$R_{2,40}$ & Diagonal-fermion quotient of two identical half-$E_6$ categories & U & \cite{ALW}\\
$R_{2,49}$ & Quadratic category $R^{\prime}(2)$ & U & \cite{MY}\\
$R_{2,53}$ & $\Vec_{C_2}\boxtimes\operatorname{Ad}(E_6)$; also $R(0,2)$ & U & \cite{MY}\\
$R_{2,69}$ & $C_2$-extension of $\Rep(H_8)$ & U & \cite{CLS}\\
$R_{2,71}$ & $\mathcal C(S_4,\omega,C_2^2,\psi)$ of Theorem~\ref{thm:s4-model} & U & \cite{GNgt}\\
$R_{2,75}$ & $E_6$ modules over $SU(2)_{10}$ & U & \cite{KO}\\
$R_{2,94}$ & Diagonal-fermion quotient of $\mathrm{Ising}\boxtimes PSU(2)_6$ & U & \cite{EGNO}\\
$R_{2,97}$ & Dual Goodman--de la Harpe--Jones $3311$ even part & U & \cite{IJMS}\\
$R_{2,98}$ & Dual even part of the Izumi $C_4$ subfactor & U & \cite{IMP}\\
$R_{2,105}$ & $\mathcal C(A_5,1,D_{10},1)$ & U & \cite{Integral}\\
$R_{2,124}$ & Schopieray: $\mathcal C(\mathfrak{so}_5,3/2)_{\mathrm{ad}}$ & S & \cite{SchNonPU}\\
$R_{2,127}$ & $\mathrm{Fib}\boxtimes\operatorname{Ad}(E_6)$ & U & \cite{KO}\\
$R_{2,128}$ & $D_6$ modules over $SU(2)_8$ & U & \cite{KO}\\
$R_{2,140}$ & Local $D_{10}$ modules over $SU(2)_{16}$ & U & \cite{KO}\\
$R_{2,143}$ & Commutative even part of the Izumi $C_5$ subfactor & U & \cite{EG}\\
$R_{2,144}$ & Principal $G_2$ level-three category & U & \cite{QGsurvey}\\
\addlinespace
$R_{3,47}$ & $\Rep(GL(3,\FF_2))\simeq\Rep(PSL(2,7))$ & U & \cite{Integral}\\
$R_{3,141}$ & $\mathcal C(A_5,1,S_3,1)$ & U & \cite{Integral}\\
$R_{3,218}$ & Inversion equivariantization of Izumi's $C_5+5$ near-group category & U & \cite{Izumi}\\
\addlinespace
$R_{4,249}$ & Sign-swap equivariantization of the equal half-$E_6$ quotient & U & \S\ref{sec:sign-swap}\\
$R_{4,322}$ & Diagonal-fermion quotient of half-$E_6\boxtimes PSU(2)_6$ & U & \S\ref{sec:mixed-fermion}\\
\addlinespace
$R_{5,17}$ & Izumi's $C_5+5$ near-group category & U & \cite{Izumi}\\
\addlinespace
$R_{6,345}$ & $\mathcal C_{4,8,4}$, with ring $R(4,6)$ & U & \cite{MY}\\
$R_{6,1712}$ & Relative Fibonacci center $Z_{\mathrm{Fib}}(TT_3)$ & U & \S\ref{sec:relative-fibonacci}\\
$R_{6,1764}$ & Asaeda--Haagerup category $\mathcal{AH}_1$ & U & \cite{GSAH}\\
\addlinespace
$R_{7,572}$ & $\Rep((C_3\times C_3)\rtimes Q_8)$ & U & \cite{EGNO}\\
$R_{7,2274}$ & Extended Haagerup category $\mathcal{EH}_1$ & U & \cite{BMPS,BPEH}\\
\end{longtable}
\endgroup
